% =============================================================================
% Chapter 3 -- The State-Vector Seam
% =============================================================================
%
% Audience prerequisite: Chapter 2 (installation and first run).
%
% Reader takeaway: after this chapter, the reader can convert a host-code
%   population array of shape (N_tau, L) to GRAIL's flat blocked form
%   (L*N_tau,) and back, knowing exactly which slot holds which population
%   at which depth.
%
% Cross-reference note: Chapter 4 defines its own \label{sec:seam} for its
%   brief seam subsection (§4.2).  This chapter uses \label{sec:two-conventions}
%   for the canonical treatment so there is no label collision in the master.
%   When a future Ch 4 copy-edit adds a \ref back here, use \ref{sec:two-conventions}.
% =============================================================================

\chapter{The State-Vector Seam}
\label{ch:statevec}

\refguide{6.5}{Tensor Conventions}

Every call to \code{resolvent\_solve} crosses one boundary.
On one side is your host code's array layout.
On the other is GRAIL's flat layout.
That boundary is the state-vector seam.
Getting it right once makes every subsequent JVP callback correct.
Getting it wrong produces a wrong answer with no error message.

% =============================================================================
\section{The Two Conventions}
\label{sec:two-conventions}
% =============================================================================

Inside your host code, atomic populations live in a 2-D array of shape
\code{(N\_tau, L)}.
Each row is a depth point; each column is an atomic level.
At the GRAIL boundary, the same data lives in a 1-D flat array of shape
\code{(L*N\_tau,)}.
The two conventions exist for different reasons.
The host code thinks per-depth.
Row~$i$ holds all level populations at optical depth~$\tau_i$.
GRAIL is dimension-agnostic.
It treats state as a flat vector.
It never inspects the physical meaning of any entry.
One packing rule connects them: level-major order.
All depths for level~1 come first, then all depths for level~2, and so on.

% =============================================================================
\section{The Pack}
\label{sec:pack}
% =============================================================================

The operation that converts host-side to GRAIL-side is
\code{populations.T.flatten()}.

Here is what that means on a small concrete array.

\grailInLabel{1}
\begin{lstlisting}[style=grailInStyle]
import torch

N_tau, L = 3, 2   # 3 depth points, 2 atomic levels

# populations[i, l] = fraction of atoms at depth i in level l.
populations = torch.tensor([
    [0.9500, 0.0500],   # depth 0: n_1, n_2
    [0.9200, 0.0800],   # depth 1
    [0.8900, 0.1100],   # depth 2
], dtype=torch.float64)   # shape (N_tau, L) = (3, 2)

print("shape:", populations.shape)
print(populations)
\end{lstlisting}

\grailOutLabel{1}
\begin{lstlisting}[style=grailOutStyle]
shape: torch.Size([3, 2])
tensor([[0.9500, 0.0500],
        [0.9200, 0.0800],
        [0.8900, 0.1100]], dtype=torch.float64)
\end{lstlisting}

The transpose swaps the axes so that levels become rows and depths become
columns.

\grailInLabel{2}
\begin{lstlisting}[style=grailInStyle]
# Step 1: transpose -> shape (L, N_tau) = (2, 3).
print("transposed:")
print(populations.T)
\end{lstlisting}

\grailOutLabel{2}
\begin{lstlisting}[style=grailOutStyle]
transposed:
tensor([[0.9500, 0.9200, 0.8900],   # level 1 at depths 0, 1, 2
        [0.0500, 0.0800, 0.1100]], dtype=torch.float64)  # level 2
\end{lstlisting}

Flattening that row-by-row produces the GRAIL-side vector.

\grailInLabel{3}
\begin{lstlisting}[style=grailInStyle]
# Step 2: flatten -> shape (L * N_tau,) = (6,).
x_flat = populations.T.flatten()
print("flat shape:", x_flat.shape)
print(x_flat)
print()
print("Slots [0:N_tau]       -> n_1 at all depths:", x_flat[:N_tau])
print("Slots [N_tau:2*N_tau] -> n_2 at all depths:", x_flat[N_tau:])
\end{lstlisting}

\grailOutLabel{3}
\begin{lstlisting}[style=grailOutStyle]
flat shape: torch.Size([6])
tensor([0.9500, 0.9200, 0.8900, 0.0500, 0.0800, 0.1100], dtype=torch.float64)

Slots [0:N_tau]       -> n_1 at all depths: tensor([0.9500, 0.9200, 0.8900], ...)
Slots [N_tau:2*N_tau] -> n_2 at all depths: tensor([0.0500, 0.0800, 0.1100], ...)
\end{lstlisting}

\grailNote{The first \code{N\_tau} entries are level-1 populations across
all depths.  The next \code{N\_tau} entries are level-2 populations across
all depths.  For $L$ levels, the general indexing rule is:}

\[
  x_\text{flat}[k \cdot N_\tau + i] = n_{k+1}(\tau_i),
  \quad k = 0, \ldots, L-1,\; i = 0, \ldots, N_\tau - 1.
\]

The indexing formula makes every slot traceable.

\grailInLabel{4}
\begin{lstlisting}[style=grailInStyle]
# Verify the indexing formula: x_flat[k * N_tau + i] = n_{k+1}(tau_i).
k, i = 1, 2   # level 2 (k=1), depth 2 (i=2)
expected = populations[i, k]          # host-side access: populations[depth, level]
actual   = x_flat[k * N_tau + i]      # GRAIL-side access
print(f"populations[{i}, {k}] = {expected.item():.4f}")
print(f"x_flat[{k}*{N_tau}+{i}] = x_flat[{k*N_tau+i}] = {actual.item():.4f}")
print("match:", torch.isclose(expected, actual).item())
\end{lstlisting}

\grailOutLabel{4}
\begin{lstlisting}[style=grailOutStyle]
populations[2, 1] = 0.1100
x_flat[1*3+2] = x_flat[5] = 0.1100
match: True
\end{lstlisting}

Pass this vector as \code{x\_star} and \code{delta\_T} to
\code{resolvent\_solve}.
Every JVP callback must accept and return a vector in the same flat layout.

% =============================================================================
\section{The Unpack}
\label{sec:unpack}
% =============================================================================

The inverse of \code{populations.T.flatten()} is
\code{flat.reshape(L,\,N\_tau).T}.
Every JVP callback begins by unpacking its input vector with this operation.

Here is the complete round-trip.

\grailInLabel{5}
\begin{lstlisting}[style=grailInStyle]
# Pack.
x_flat = populations.T.flatten()   # (L*N_tau,)

# Unpack.
populations_recovered = x_flat.reshape(L, N_tau).T   # (N_tau, L)

print("recovered:")
print(populations_recovered)
print("round-trip exact:", torch.allclose(populations, populations_recovered))
\end{lstlisting}

\grailOutLabel{5}
\begin{lstlisting}[style=grailOutStyle]
recovered:
tensor([[0.9500, 0.0500],
        [0.9200, 0.0800],
        [0.8900, 0.1100]], dtype=torch.float64)
round-trip exact: True
\end{lstlisting}

The round-trip is exact.
Every JVP callback packs its output and unpacks its input.
If either operation uses the wrong convention, GMRES produces a wrong
$\delta x^*$.
No error message appears.
Run this round-trip check as a unit test before calling
\code{resolvent\_solve}.

% =============================================================================
\section{Why \texttt{.T.flatten()} and Not \texttt{.flatten()}}
\label{sec:why-transpose}
% =============================================================================

\code{populations.flatten()} (without the transpose) produces a different
order.

\grailInLabel{6}
\begin{lstlisting}[style=grailInStyle]
# WRONG: no transpose -- produces depth-major (interleaved) order.
x_wrong = populations.flatten()
print("wrong:", x_wrong)
print("entry 0:", x_wrong[0], "<- n_1(tau_0)")
print("entry 1:", x_wrong[1], "<- n_2(tau_0)  [level 2, not depth 1]")
\end{lstlisting}

\grailOutLabel{6}
\begin{lstlisting}[style=grailOutStyle]
wrong: tensor([0.9500, 0.0500, 0.9200, 0.0800, 0.8900, 0.1100], ...)
entry 0: 0.9500 <- n_1(tau_0)
entry 1: 0.0500 <- n_2(tau_0)  [level 2, not depth 1]
\end{lstlisting}

\code{populations.flatten()} interleaves levels and depths: the ordering is
$n_1(\tau_0),\, n_2(\tau_0),\, n_1(\tau_1),\, n_2(\tau_1), \ldots$.
GRAIL expects the opposite: one level across all depths, then the next.

\grailInLabel{7}
\begin{lstlisting}[style=grailInStyle]
# CORRECT: transpose first, then flatten -- produces level-major order.
x_right = populations.T.flatten()
print("right:", x_right)
print("entry 0:", x_right[0], "<- n_1(tau_0)")
print("entry 1:", x_right[1], "<- n_1(tau_1)  [same level, next depth]")
print("entry 3:", x_right[3], "<- n_2(tau_0)  [level 2 starts at index N_tau]")
\end{lstlisting}

\grailOutLabel{7}
\begin{lstlisting}[style=grailOutStyle]
right: tensor([0.9500, 0.9200, 0.8900, 0.0500, 0.0800, 0.1100], ...)
entry 0: 0.9500 <- n_1(tau_0)
entry 1: 0.9200 <- n_1(tau_1)  [same level, next depth]
entry 3: 0.0500 <- n_2(tau_0)  [level 2 starts at index N_tau]
\end{lstlisting}

\begin{advisorybox}
\textbf{Rule:} Use \code{populations.T.flatten()} to pack and
\code{flat.reshape(L, N\_tau).T} to unpack.
Both operations must appear, consistently, in every JVP callback.
The interleaved \code{.flatten()} order is the most common layout bug.
It passes all shape checks.
It produces a silently wrong $\delta x^*$.
\end{advisorybox}

% =============================================================================
\section{Multi-Level State Vectors}
\label{sec:multiline}
% =============================================================================

The same convention applies for any $L$.
The 3-level toy hydrogen atmosphere has $L = 3$ levels and $K = 2$ lines:
Ly-$\alpha$ and H-$\alpha$.

\grailInLabel{8}
\begin{lstlisting}[style=grailInStyle]
# 3-level example: N_tau = 4, L = 3.
N_tau, L = 4, 3

pops3 = torch.rand(N_tau, L, dtype=torch.float64)
# Normalise rows so populations sum to 1.
pops3 = pops3 / pops3.sum(dim=1, keepdim=True)

x_flat3 = pops3.T.flatten()   # (L*N_tau,) = (12,)

print("pops3 shape:", pops3.shape)
print("flat  shape:", x_flat3.shape)
print()
print("Slots [0:N_tau]         = n_1 over depths:", x_flat3[:N_tau])
print("Slots [N_tau:2*N_tau]   = n_2 over depths:", x_flat3[N_tau:2*N_tau])
print("Slots [2*N_tau:3*N_tau] = n_3 over depths:", x_flat3[2*N_tau:])
\end{lstlisting}

\grailOutLabel{8}
\begin{lstlisting}[style=grailOutStyle]
pops3 shape: torch.Size([4, 3])
flat  shape: torch.Size([12])

Slots [0:N_tau]         = n_1 over depths: tensor([...], dtype=torch.float64)
Slots [N_tau:2*N_tau]   = n_2 over depths: tensor([...], dtype=torch.float64)
Slots [2*N_tau:3*N_tau] = n_3 over depths: tensor([...], dtype=torch.float64)
\end{lstlisting}

The indexing formula $x_\text{flat}[k \cdot N_\tau + i] = n_{k+1}(\tau_i)$
holds for any $L$.
GRAIL receives a vector of length $L \cdot N_\tau$.
It never inspects $L$ directly.
You carry $L$ in your callback closures.

The mean radiation field $\bar J$ per spectral line uses a parallel
convention.
For a $K$-line atmosphere, the host-side array is \code{J\_bar\_lines} of
shape \code{(N\_tau, K)}.

\grailInLabel{9}
\begin{lstlisting}[style=grailInStyle]
# J_bar_lines convention: (N_tau, K) -> flat (K*N_tau,).
N_tau, K = 4, 2   # 4 depths, 2 lines (Ly-alpha and H-alpha)

J_bar_lines = torch.rand(N_tau, K, dtype=torch.float64)

J_flat = J_bar_lines.T.flatten()   # (K*N_tau,)

print("J_bar_lines shape:", J_bar_lines.shape)
print("J_flat shape:     ", J_flat.shape)
print()
print("Slots [0:N_tau]       = J_bar_Lya over depths:", J_flat[:N_tau])
print("Slots [N_tau:2*N_tau] = J_bar_Ha  over depths:", J_flat[N_tau:])
\end{lstlisting}

\grailOutLabel{9}
\begin{lstlisting}[style=grailOutStyle]
J_bar_lines shape: torch.Size([4, 2])
J_flat shape:      torch.Size([8])

Slots [0:N_tau]       = J_bar_Lya over depths: tensor([...], dtype=torch.float64)
Slots [N_tau:2*N_tau] = J_bar_Ha  over depths: tensor([...], dtype=torch.float64)
\end{lstlisting}

\grailNote{The $\bar J$ flat layout is used inside host JVP callbacks to
pack and unpack the radiation-field tangent, not in the GRAIL state vector
itself.  The GRAIL state vector contains populations only.}

In the 3-level host JVP, the callback receives a populations tangent of
shape \code{(3*N\_tau,)}.
It produces $\bar J$ tangents internally and returns a flat output of shape
\code{(3*N\_tau,)}.
The $\bar J$ flat layout appears only in the intermediate step.
Chapter~\ref{ch:wrapping} shows the full wiring.

% =============================================================================
\section{Quick Reference}
\label{sec:seam-ref}
% =============================================================================

The table below collects every layout operation in one place.

\begin{argtable}
\code{populations}          & \code{Tensor} & \code{(N\_tau, L)} &
  Host-side array.  Row~$i$ holds all level populations at depth~$\tau_i$.
  Dtype: float64. \\
\code{populations.T.flatten()} & \code{Tensor} & \code{(L*N\_tau,)} &
  GRAIL-side flat vector.  Entries \code{[k*N\_tau : (k+1)*N\_tau]} hold
  $n_{k+1}(\tau)$ across all depths.  Pass as \code{x\_star} and
  \code{delta\_T}. \\
\code{flat.reshape(L,\,N\_tau).T} & \code{Tensor} & \code{(N\_tau, L)} &
  Inverse of the pack.  Use at the start of every JVP callback to
  recover the 2-D layout. \\
\code{J\_bar\_lines}        & \code{Tensor} & \code{(N\_tau, K)} &
  Host-side radiation field for $K$ spectral lines.  Row~$i$ holds the
  $K$ line-averaged intensities at depth~$\tau_i$.  Dtype: float64. \\
\code{J\_bar\_lines.T.flatten()} & \code{Tensor} & \code{(K*N\_tau,)} &
  Flat layout of the radiation field.  Entries \code{[k*N\_tau :
  (k+1)*N\_tau]} hold $\bar J_k(\tau)$ across all depths.  Used inside
  host JVP callbacks; not part of the GRAIL state vector. \\
\end{argtable}

\grailInLabel{10}
\begin{lstlisting}[style=grailInStyle]
# Minimal correctness check -- run before every new integration.
def check_seam(populations: torch.Tensor) -> bool:
    """Return True if the pack/unpack round-trip is exact."""
    N_tau, L = populations.shape
    x_flat = populations.T.flatten()
    recovered = x_flat.reshape(L, N_tau).T
    return torch.allclose(populations, recovered)

# Example: 2-level slab, 40 depth points.
N_tau, L = 40, 2
pops_test = torch.rand(N_tau, L, dtype=torch.float64)
pops_test = pops_test / pops_test.sum(dim=1, keepdim=True)

print("seam check:", check_seam(pops_test))
\end{lstlisting}

\grailOutLabel{10}
\begin{lstlisting}[style=grailOutStyle]
seam check: True
\end{lstlisting}

\grailNote{Add \code{assert check\_seam(populations)} as the first line
of your integration test.  A failure here means the JVP callbacks will
operate in the wrong coordinate system for every GMRES iteration.}

Chapter~\ref{ch:wrapping} builds the full host JVP and emulator JVP on top
of this convention.
Each callback begins with \code{v\_pops = v.reshape(L, n\_tau).T} and ends
with \code{return result.T.flatten()}.
The seam is the only thing that makes those two lines safe.
